\documentclass[../main]{subfiles}
\begin{document}

\section{Mantenimiento Predictivo}

\imgfull{images/09/logo2}{antenimiento predictivo}

\info{Mantenimiento predictivo}{El mantenimiento predictivo es un tipo de mantenimiento que relaciona una variable física con el desgaste o estado de una máquina. El mantenimiento predictivo se basa en la medición, seguimiento y monitoreo de parámetros y condiciones operativas de un equipo o instalación.}

El mantenimiento predictivo incluye muchas actividades que tienen que ver con herramientas técnicas y científicas para el estudio de conceptos físico, químicos con el fin de lograr la máxima fiabilidad en los equipos.



\subsection{Análisis de vibraciones}

\img{images/09/logo}{Análisis de vibraciones en acoplamiento.}

El análisis de vibraciones es la principal técnica para supervisar y diagnosticar la maquinaria rotativa e implantar un plan de mantenimiento predictivo.
El análisis de vibraciones se aplica con eficacia desde hace más de 30 años a la supervisión y diagnóstico de fallos mecánicos en máquinas rotativas. Inicialmente, se emplearon equipos analógicos para la medida de la vibración en banda ancha, lo que hacía imposible el diagnóstico fiable de fallos en rodamientos y engranajes.

Más tarde, se incorporaron filtros sintonizables a la electrónica analógica, lo que incrementó enormemente la capacidad de diagnóstico, pero sin poder tratar la información de forma masiva. Desde 1984, se comenzaron a emplear equipos digitales con FFT en tiempo real y capacidad de almacenamiento (analizadores-colectores) y tratamiento en software para PC.

Hoy día nadie pone en duda la capacidad del análisis de vibraciones en máquinas rotativas, que incluso permite el diagnóstico de algunos problemas en máquinas eléctricas.

\subsubsection{Fallos detectables}
Mediante el análisis de vibraciones aplicado a la maquinaria rotativa se pueden diagnosticar con precisión problemas de:
\begin{itemize}
\item Desequilibrio y Desalineación
\item Holguras y Roces
\item Ejes doblados y Poleas excéntricas
\item Rodamientos y Engranajes
\item Fallos de origen eléctrico
\end{itemize}

\subsubsection{Maquinaria crítica monitorizable} 
La maquinaria crítica susceptible de ser monitorizada en las plantas industriales es la siguiente:
\begin{itemize}
\item Turbinas de vapor y de gas
\item Bombas centrífugas
\item Ventiladores
\item Motores eléctricos
\item Compresores rotativos, de tornillo y alternativos
\item Agitadores, mezcladoras...
\item Molinos y hornos rotativos
\item Cajas reductoras Centrífugas
\item Torres de refrigeración
\item Motores diesel y generadores de equipos electrógenos
\end{itemize}



\img{images/09/analisisvibraciones}{Análisis de Vibraciones.}


\subsection{Análisis de radiación infraroja}


\img{images/09/termografia}{Termografía en elementos eléctricos.}

La termografía infrarroja es la ciencia de adquisición y análisis de información térmica a partir de dispositivos de toma de imágenes sin contacto. Es una técnica de mantenimiento predictivo que puede ser usado para el monitoreo de la condición de máquinas, estructuras y sistemas.

Esta usa instrumentos diseñados para la detección de emisión de energía infrarroja, para determinar la condición ya sea de sistemas electromecánicos, procesos de producción (Hornos, calderas, etc.), edificios o casas y hasta de humanos.
Se pueden detectar anomalías, por ejemplo áreas que están más calientes o más frías de lo que deberían estar, un experimentado inspector puede localizar y definir problemas incipientes dentro de una planta.

\subsubsection{¿Qué es la radiación infrarroja?}

La radiación infrarroja, o radiación IR es un tipo de radiación electromagnética, de mayor longitud de onda que la luz visible, pero menor que la de las microondas. Por ello, tiene menor frecuencia que la luz visible y mayor que las microondas. Su rango de longitudes de onda va desde unos 0,7 hasta los 1000 micrómetros. La radiación infrarroja es emitida por cualquier cuerpo cuya temperatura sea mayor que 0 Kelvin, es decir -273,15$^oC$ (cero absoluto).

\info{Termografía}{En Termografía se utiliza la más moderna tecnología en cámaras infrarrojas, algunas aplicaciones que podemos mencionar:
\begin{itemize}
\item Detección de fiebre o síntomas de enfermedades.
\item Detección de Cáncer y otras aplicaciones médicas
\item Inspección de Hornos
\item Inspección de aviones para detección de humedad en fuselajes. 
\end{itemize}
}
\img{images/09/termografia2}{Termografía.}

\newpage

\actividad{Leer, interpretar y contestar}
\begin{enumerate}
    \item ¿En que consiste el mantenimiento predictivo?
    \item ¿Qué actividades comprende el mantenimiento predictivo?
    \item ¿En qué consiste el análisis de vibraciones?
    \item ¿Qué fallos son posibles detectar por análisis de vibraciones y que maquinaria?
    \item ¿Qué es la termografía y cuál es su uso en mantenimiento?
\end{enumerate}
\actividad{Investigar}
\begin{enumerate}
     \item ¿Qué es la radiación la radiación electromagnética infrarroja?
     \item ¿Qué es la espectroscopia infrarroja y que estudia?
     \item ¿Qué es la radiancia espectral?
     \item ¿Qué es el espectro electromagnético?
     \item ¿Qué es la reflectografía infrarroja?
    
\end{enumerate}
\end{document}
